% ---------------------------------------------------------------------------
%  QML Lab research project 09: Quantum Transfer Learning
%  Build:  latexmk -pdf main.tex   (needs qmlproject.sty, harvardqml.bst, references.bib)
% ---------------------------------------------------------------------------
\documentclass[11pt,a4paper]{article}
\usepackage{qmlproject}
\qmlheader{Project 09: Quantum Transfer Learning}

\hypersetup{pdftitle={QML Lab project 09: Quantum Transfer Learning}}

\begin{document}

\begin{center}
  {\small\textcolor{black!55}{QML Lab \textbullet\ research project 09}}\\[0.5em]
  {\LARGE\bfseries Quantum Transfer Learning}\\[0.6em]
  {\large A pretrained ResNet as feature extractor and a small quantum circuit as the classifier head.}\\[0.6em]
  {\small Living document \textbullet\ last updated \today}
\end{center}

\begin{projectcard}
\small
\begin{tabularx}{\linewidth}{@{}>{\bfseries}l X >{\bfseries}l X@{}}
Quadrant  & \CQ\ (classical data, quantum processor) & Track & Hybrid \\
Difficulty & Easy & Tools & PennyLane, PyTorch, torchvision \\
Builds on & Module 7 of the QML Lab & Workspace & \path{qml_projects/quantum-transfer-learning/} \\
\end{tabularx}
\end{projectcard}

\paragraph{Summary.}
Transfer learning reuses a network trained on a large data set and retrains only a small final part
for a new task. Here the reused part is a pretrained ResNet-18 and the retrained part is a
``dressed'' quantum circuit: a classical layer, a 4-qubit variational circuit and another classical
layer. The project measures whether the quantum head matches classical heads of the same size, and
how the result depends on qubit count and on the data set.

% ===========================================================================
\section{Background and key concepts}
% ===========================================================================
\subsection{Transfer learning}
Deep convolutional networks such as ResNet-18 \citep{he2016} learn general visual features on
ImageNet. Removing the final layer leaves a frozen feature extractor that maps an image to a vector
$z\in\R^{512}$. Training only a new head on $z$ works well even with a few hundred labelled images.

\subsection{Dressed quantum circuits}
\citet{mari2020} proposed classical-to-quantum transfer learning with a head of the form
\begin{equation}
  \hat y=W_{\text{out}}\;\big\langle Z\big\rangle_\theta\!\Big(\tfrac{\pi}{2}\tanh\big(W_{\text{in}}z\big)\Big),
\end{equation}
where $W_{\text{in}}\in\R^{n\times512}$ reduces the features to $n$ angles, the variational circuit
returns one $\langle Z\rangle$ per qubit, and $W_{\text{out}}$ maps them to class scores. Only the head is
trained, with PyTorch \citep{paszke2019}.

\subsection{What a fair comparison looks like}
The backbone does most of the work, so the comparison must hold everything else fixed: the same
frozen features, the same training budget, and classical heads with a similar number of parameters
(for example $512\to n\to2$ with a nonlinearity in place of the circuit).

% ===========================================================================
\section{Research questions}
% ===========================================================================
The project is open-ended: these questions are the starting point, not the finish line. Add new
ones to the README and the journal as they come up.

\begin{rqlist}
  \item Does the quantum head match a classical head with the same number of parameters?
  \item How does accuracy change with 2, 4 and 6 qubits and with circuit depth?
  \item How much does the $512\to n$ reduction layer matter compared with the circuit itself?
\end{rqlist}

% ===========================================================================
\section{Method}
% ===========================================================================
\subsection{Data}
The two-class ants/bees image set used by \citet{mari2020} (a few hundred images), then a two-class
subset of CIFAR-10. Standard ImageNet normalisation and light augmentation; fixed train/validation
splits per seed.

\subsection{Models}
Frozen ResNet-18 from torchvision; dressed quantum circuit with $n=4$ qubits, Hadamard layer,
angle encoding and depth 6 of $R_Y$ rotations and CNOT chains; cross-entropy, Adam, 15 epochs,
batch size 4.

\subsection{Baselines}
Linear head $512\to2$; MLP head $512\to n\to2$ with $\tanh$ (same parameter count as the dressed
circuit apart from the circuit angles); fine-tuning the whole network as an upper bound.

\subsection{Metrics}
\begin{itemize}
  \item Validation accuracy (mean $\pm$ s.d. over seeds).
  \item Number of trainable head parameters; training time per epoch.
\end{itemize}

% ===========================================================================
\section{Planned experiments}
% ===========================================================================
Each experiment becomes one \texttt{[[experiments]]} entry in the workspace, so its runs, parameters
and metrics are stored and compared in the portal.

\begin{table}[h]
\centering\small
\begin{tabularx}{\linewidth}{@{}l>{\raggedright}X>{\raggedright}X>{\raggedright\arraybackslash}X@{}}
\toprule
ID & Experiment & What varies & Main metrics \\
\midrule
  E1 & Quantum vs classical head & head type, $n=4$ & accuracy \\
  E2 & Qubit sweep & $n\in\{2,4,6\}$ & accuracy \\
  E3 & Circuit depth & depth $1..8$ & accuracy, time \\
  E4 & Noise & noisy simulator at inference & accuracy drop \\
  E5 & Second data set & CIFAR-10 pair & accuracy \\
\bottomrule
\end{tabularx}
\caption{Experiment plan. Results go into the journal and the results log (\cref{sec:log}).}
\end{table}

% ===========================================================================
\section{Running it in the QML Lab}
% ===========================================================================
The workspace lives in \path{qml_projects/quantum-transfer-learning/}. Scripts import three helpers from
\path{qml_lab/plots.py}:
\begin{itemize}
  \item \texttt{params()} reads the parameters from the portal's run form (defaults when run by hand);
  \item \texttt{save()} stores figures with the run;
  \item \texttt{log\_metrics()} records numbers that appear in the run table.
\end{itemize}

\begin{lstlisting}
qml_projects/quantum-transfer-learning/
  project.toml          metadata + experiments
  README.md             goal, open questions, ideas
  experiments/          runnable scripts
  docs/main.tex         this document
\end{lstlisting}

First experiment: This experiment is \emph{planned}: write the script, add the block below to \texttt{project.toml}, then run \texttt{python manage.py sync\_projects}.

\begin{lstlisting}[language={}]
[[experiments]]
slug = "ants-bees"
title = "Dressed quantum head on ResNet-18 features"
script = "experiments/e01_ants_bees.py"
figure = "head_comparison.png"

[experiments.params]
n_qubits = 4
q_depth = 6
epochs = 15
lr = 0.0004
batch_size = 4
seed = 0
\end{lstlisting}

\begin{lstlisting}[language=bash]
# from the portal folder: pick up changes to project.toml
python manage.py sync_projects
cd qml_projects/quantum-transfer-learning
# run by hand with the default parameters
PYTHONPATH=../../qml_lab python experiments/e01_ants_bees.py
\end{lstlisting}

% ===========================================================================
\section{Quality bar}
% ===========================================================================
\begin{itemize}
  \item A classical baseline with a similar number of trainable parameters, tuned with the same care \citep{bowles2024}.
  \item At least five random seeds per configuration, reported as mean $\pm$ standard deviation.
  \item How results change with the number of qubits.
  \item At least one run under a realistic noise model.
  \item Compare heads only on identical frozen features; report head parameter counts.
\end{itemize}

% ===========================================================================
\section{Open directions}
% ===========================================================================
\begin{itemize}
  \item Quantum-to-classical transfer: a pretrained quantum feature extractor with a classical head.
  \item Inference of the trained head on real hardware (Project 16).
  \item Other backbones (MobileNet, vision transformers).
\end{itemize}

% ===========================================================================
\section{Risks and limitations}
% ===========================================================================
\begin{itemize}
  \item The backbone dominates; small differences between heads may be noise.
  \item Tiny data sets give high variance; use several seeds.
\end{itemize}

% ===========================================================================
\section{Results log}
\label{sec:log}
% ===========================================================================
Summarise findings here as they become stable; day-to-day notes belong in the portal journal.

\begin{table}[h]
\centering\small
\begin{tabularx}{\linewidth}{@{}l l >{\raggedright}X >{\raggedright\arraybackslash}X@{}}
\toprule
Date & Run & Finding & Next step \\
\midrule
 & & & \\
 & & & \\
 & & & \\
\bottomrule
\end{tabularx}
\end{table}

\bibliographystyle{harvardqml}
\bibliography{references}

\end{document}
